\section*{Chapitre \ref{ch:fm:compliance-and-market-abuse} --- Conformité et abus de marché}

\begin{solution}{exo:fm:compliance-and-market-abuse:1}
La règle de blackout~: la société a négocié le titre dans les deux jours qui précèdent la demande. Elle empêche les salariés d'opérer sur la base de la connaissance des ordres propres de la société,
en les devançant ou en les accompagnant.
\end{solution}

\begin{solution}{exo:fm:compliance-and-market-abuse:2}
$100/11\,600=0.86~\%$.
\end{solution}

\begin{solution}{exo:fm:compliance-and-market-abuse:3}
La \omterm{def:fm:compliance-and-market-abuse:lists}{liste d'interdiction} bloque les opérations et son existence est connue des salariés~; la \omterm{def:fm:compliance-and-market-abuse:lists}{liste de surveillance}, réservée à la conformité, recense les noms pour lesquels une partie de la société détient
une \omterm{def:fm:compliance-and-market-abuse:abuse}{information privilégiée}. La publier reviendrait à transmettre elle-même cette information.
\end{solution}

\begin{solution}{exo:fm:compliance-and-market-abuse:4}
Les 580 alertes dépassent la capacité de 100, et les analystes les lisent dans l'ordre d'arrivée~: en moyenne 17~\% des 89 alertes véritables sont lues,
$89\times100/580=15.3$.
\end{solution}

\begin{solution}{exo:fm:compliance-and-market-abuse:5}
Les teneurs de marché légitimes annulent presque tous leurs ordres, davantage que les auteurs de spoofing, si bien qu'un ratio élevé placerait les teneurs de marché en tête~; leur signature
tient à l'écart entre les taux d'exécution de leurs gros et de leurs petits ordres, et au moment de leurs annulations.
\end{solution}

\begin{solution}{exo:fm:compliance-and-market-abuse:6}
Jours 40 à 74~: du \omterm{def:fm:compliance-and-market-abuse:barrier}{franchissement de la muraille} jusqu'à la purge du jour 75. L'\omterm{def:fm:compliance-and-market-abuse:pad}{autorisation préalable} ne couvre que le compte propre de la personne ayant franchi la muraille~; l'information pourrait
atteindre le desk de trading ou d'autres personnes, ce que seule une surveillance de toutes les transactions sur ce titre permettrait de détecter.
\end{solution}

\begin{solution}{exo:fm:compliance-and-market-abuse:7}
Dans l'ordre d'arrivée, le meilleur taux est 2.5~\% et en trouve 81~; le taux par défaut de 5~\% en trouve 46. Dans l'ordre des scores, le taux par défaut en trouve aussi 81~: examiner en premier
les scores les plus élevés rend inoffensif un seuil bas.
\end{solution}

\begin{solution}{exo:fm:compliance-and-market-abuse:8}
Davantage d'alertes avec les mêmes analystes signifie une part lue plus faible~: au-delà du taux de capacité, abaisser le seuil réduit le nombre de cas véritables trouvés dans
l'ordre d'arrivée (de 65 pour 0.8~\% à 15 pour 5~\% dans le chapitre). Le critère est le nombre de cas véritables trouvés, non celui des alertes générées.
\end{solution}

\begin{solution}{pb:fm:compliance-and-market-abuse:1}
\begin{enumerate}
\item Voir la \cref{def:fm:compliance-and-market-abuse:abuse}.
\item Des dispositifs, systèmes et procédures efficaces pour détecter et déclarer les ordres et transactions suspects, et une notification sans délai en cas de
soupçon raisonnable.
\item Voir l'\cref{dat:fm:compliance-and-market-abuse:us}.
\item Elle n'a pas correctement appliqué l'exigence de 2016, a mis dix-huit mois à évaluer ses risques et n'a pu surveiller efficacement certains
types d'abus~: une amende de 12.6~millions de livres sterling.
\item Voir la \cref{def:fm:compliance-and-market-abuse:stor}.
\item Voir la \cref{prop:fm:compliance-and-market-abuse:queue}.
\item 0.8~\% et 65 cas trouvés avec 50 heures~; 2.5~\% et 81 avec 150.
\item 580 alertes, dont 89 véritables, 15 trouvées et 85 manquées.
\item Au taux par défaut, il en trouve 66 avec 50 heures et 81 avec 150~: c'est l'ordre de revue, et non le seuil, qui décide de ce qui est lu.
\item Une alerte non lue n'est pas trouvée~: le rappel de la fonction est égal au rappel du détecteur multiplié par la part des alertes véritables examinées.
\item Voir les \cref{def:fm:compliance-and-market-abuse:pad,def:fm:compliance-and-market-abuse:lists}.
\item 964 demandes approuvées~; refusées~: 202 pour négociation récente de la société, 23 pour \omterm{def:fm:compliance-and-market-abuse:lists}{liste d'interdiction}, 15 pour initiés, 6 pour période de détention et 5 pour \omterm{def:fm:compliance-and-market-abuse:lists}{liste de surveillance}.
\item Voir la \cref{def:fm:compliance-and-market-abuse:barrier}.
\item Du franchissement jusqu'à la purge~; toute transaction personnelle sur les instruments de l'émetteur est refusée entre-temps.
\item Une alerte clôturée sans avoir été lue est un cas non trouvé, et une trace du choix de la société de ne pas l'examiner.
\item Afin que les alertes de chaque détecteur remplissent sa part des heures des analystes, revues dans l'ordre des scores, et ajustées au taux mesuré d'alertes véritables.
\item Associer chaque produit à l'abus qu'il permet, vérifier qu'un détecteur le couvre avant la première transaction, et tester à l'aide de cas injectés.
\item Injecter chaque année des cas connus dans les données et mesurer combien la chaîne complète, détecteurs et file d'attente, en trouve.
\item Le taux d'alerte dont les alertes remplissent juste la capacité, $q^*=C/N$ (0.86~\%, meilleur point de la grille 0.8~\%, 65 cas trouvés)~; au taux par défaut de 5~\%, la file d'attente manque
85~\% des cas.
\item Dimensionner les seuils selon l'équipe et examiner les alertes dans l'ordre des scores, avec une couverture vérifiée pour chaque nouvelle activité et testée avec des cas injectés~;
soumettre chaque transaction personnelle à une \omterm{def:fm:compliance-and-market-abuse:pad}{autorisation préalable} au regard des listes datées, et consigner chaque \omterm{def:fm:compliance-and-market-abuse:barrier}{franchissement de la muraille}.
\end{enumerate}
\end{solution}

\begin{solution}{iq:fm:compliance-and-market-abuse:1}
Information précise, non publique et sensible au cours d'un émetteur ou d'un instrument~: ne pas négocier ni la transmettre, et prévenir aussitôt la conformité, qui
placera le titre sur une \omterm{def:fm:compliance-and-market-abuse:lists}{liste d'interdiction} ou de surveillance et inscrira la personne sur la \omterm{def:fm:compliance-and-market-abuse:barrier}{liste d'initiés}.
\iqlookfor{la définition et la remontée à la conformité.}
\end{solution}

\begin{solution}{iq:fm:compliance-and-market-abuse:2}
Pour éviter de négocier sur les ordres de la société ou sur une \omterm{def:fm:compliance-and-market-abuse:abuse}{information privilégiée} détenue ailleurs dans la société, et laisser une trace attestant le contrôle effectué.
\iqlookfor{front-running, \omterm{def:fm:compliance-and-market-abuse:abuse}{information privilégiée} et trace.}
\end{solution}

\begin{solution}{iq:fm:compliance-and-market-abuse:3}
Des entrées indiquant l'instrument, la liste, les horodatages de début et de fin et le motif, en ajout seul, les corrections étant saisies comme de nouvelles entrées~; les requêtes se font à un instant donné.
\iqlookfor{bitemporel, historique en ajout seul.}
\end{solution}

\begin{solution}{iq:fm:compliance-and-market-abuse:4}
Cela dépend de la capacité~: à 5~\% d'une population importante, les alertes peuvent largement dépasser ce que les analystes peuvent lire, et les cas trouvés sont égaux au rappel multiplié par la
part lue. Demander la précision et la file d'attente.
\iqlookfor{le rappel n'est pas le résultat de la fonction.}
\end{solution}

\begin{solution}{iq:fm:compliance-and-market-abuse:5}
Le franchissement contrôlé d'une \omterm{def:fm:compliance-and-market-abuse:abuse}{information privilégiée} à travers une \omterm{def:fm:compliance-and-market-abuse:barrier}{barrière à l'information}, avec consentement~; la trace de qui, quand, quoi, du consentement, de l'inscription sur la \omterm{def:fm:compliance-and-market-abuse:barrier}{liste d'initiés},
et de la date de purge.
\iqlookfor{consentement, traces et purge.}
\end{solution}

\begin{solution}{iq:fm:compliance-and-market-abuse:6}
Donner les heures telles que l'heure suivante trouve le plus de cas véritables~: égaliser la précision de l'alerte marginale entre détecteurs, en examinant chacun dans l'ordre
des scores.
\iqlookfor{précision marginale égalisée.}
\end{solution}
